\chapter{2011年湖南卷（文）}

\section{单选题}

\begin{problem}
设全集 \(U = M \cup N = \{1, 2, 3, 4, 5\}\)，\(M \cap \complement_U N = \{2, 4\}\)，则 \(N =\)（\quad）

\choices
    {\(\{1, 2, 3\}\)}
    {\(\{1,3,5\}\)}
    {\(\{1,4,5\}\)}
    {\(\{2,3,4\}\)}
\end{problem}

\begin{answer}
B
\end{answer}

\begin{solution}
因为 \(U=M\cup N\)，所以 \(\complement_U N\subseteq M\)，从而
\(M\cap\complement_U N=\complement_U N=\{2,4\}\). 因此
\(N=\complement_U\{2,4\}=\{1,3,5\}\)，故选 B.
\end{solution}

\begin{problem}
若 \(a\)，\(b\in\R\)，\(\i\) 为虚数单位，且 \((a + \i)\i = b + \i\)，则（\quad）

\choices
    {\(a = 1\)，\(b = 1\)}
    {\(a = -1\)，\(b = 1\)}
    {\(a = 1\)，\(b = -1\)}
    {\(a = -1\)，\(b = -1\)}
\end{problem}

\begin{answer}
C
\end{answer}

\begin{solution}
由 \(\i^2=-1\)，得
\((a+\i)\i=a\i-1\). 与 \(b+\i\) 比较实部和虚部，得到
\(b=-1\)，\(a=1\)，故选 C.
\end{solution}

\begin{problem}
“\(x>1\)”是“\(|x|>1\)”的（\quad）

\choices
    {充分不必要条件}
    {必要不充分条件}
    {充分必要条件}
    {既不充分又不必要条件}
\end{problem}

\begin{answer}
A
\end{answer}

\begin{solution}
若 \(x>1\)，则 \(|x|=x>1\)，所以 \(x>1\) 是 \(|x|>1\) 的充分条件.
但 \(x=-2\) 时 \(|x|>1\) 成立而 \(x>1\) 不成立，所以它不是必要条件，故选 A.
\end{solution}

\begin{problem}
如图是某几何体的三视图，则该几何体的体积为

\bitmapfigure[width=0.6\linewidth]{img/2011/hunan_liberal/q4_fig1.png}

\choices
    {\(9\pi + 42\)}
    {\(36\pi + 18\)}
    {\(\frac{9}{2}\pi + 12\)}
    {\(\frac{9}{2}\pi + 18\)}
\end{problem}

\begin{answer}
D
\end{answer}

\begin{solution}
由三视图可知，该几何体由一个底面为边长 \(3\)、高为 \(2\) 的长方体和一个半径为 \(\frac32\) 的球组成，二者只有边界相接. 因此体积为
\(3\times3\times2+\frac43\pi\left(\frac32\right)^3=18+\frac92\pi\)，故选 D.
\end{solution}

\begin{problem}
通过随机询问 \(110\) 名性别不同的大学生是否爱好某项运动，得到如下的列联表：
\begin{center}
\begin{tabular}{c|ccc}
 & 男 & 女 & 总计 \\
\hline
爱好 & \(40\) & \(20\) & \(60\) \\
不爱好 & \(20\) & \(30\) & \(50\) \\
\hline
总计 & \(60\) & \(50\) & \(110\)
\end{tabular}
\end{center}
由 \(K^2 = \frac{n(ad-bc)^2}{(a + b)(c + d)(a + c)(b + d)}\) 算得，\(K^2 = \frac {110 \times (40 \times 30 - 20 \times 20) ^{2}}{60 \times 50 \times 60 \times 50} \approx 7.8\). 附表：
\begin{center}
\begin{tabular}{c|ccc}
\(P(K^2 \geqslant k)\) & \(0.050\) & \(0.010\) & \(0.001\) \\
\hline
\(k\) & \(3.841\) & \(6.635\) & \(10.828\)
\end{tabular}
\end{center}
参照附表，得到的正确结论是（\quad）

\choices
    {有 \(99\%\) 以上的把握认为“爱好该项运动与性别有关”}
    {有 \(99\%\) 以上的把握认为“爱好该项运动与性别无关”}
    {在犯错误的概率不超过 0.1\% 的前提下，认为“爱好该项运动与性别有关”}
    {在犯错误的概率不超过 0.1\% 的前提下，认为“爱好该项运动与性别无关”}
\end{problem}

\begin{answer}
A
\end{answer}

\begin{solution}
由题给数据，
\[
K^2=\frac{110(40\times30-20\times20)^2}{60\times50\times60\times50}\approx7.8
\]
因为 \(7.8>6.635\)，所以对应的概率小于 \(0.010\)，即在犯错误的概率不超过 \(1\%\) 的前提下，可以认为爱好该项运动与性别有关，故选 A.
\end{solution}

\begin{problem}
设双曲线 \(\frac{x^2}{a^2} - \frac{y^2}{9} = 1 \) （\(a > 0\)）的渐近线方程为 \(3x \pm 2y = 0\)，则 \(a\) 的值为（\quad）

\choices
    {4}
    {3}
    {2}
    {1}
\end{problem}

\begin{answer}
C
\end{answer}

\begin{solution}
双曲线
\(\frac{x^2}{a^2}-\frac{y^2}{9}=1\) 的渐近线为
\(y=\pm\frac{3}{a}x\). 由题意，
\(\frac{3}{a}=\frac32\)，且 \(a>0\)，所以 \(a=2\)，故选 C.
\end{solution}

\begin{problem}
曲线 \( y = \frac{\sin x}{\sin x + \cos x} - \frac{1}{2} \) 在点 \( M\left(\frac{\pi}{4}, 0\right) \) 处的切线的斜率为（\quad）

\choices
    {\(-\frac{1}{2}\)}
    {\( \frac{1}{2} \)}
    {\( -\frac{\sqrt{2}}{2} \)}
    {\( \frac{\sqrt{2}}{2} \)}
\end{problem}

\begin{answer}
B
\end{answer}

\begin{solution}
在定义域内，
\[
\left(\frac{\sin x}{\sin x+\cos x}\right)'
=\frac{\cos x(\sin x+\cos x)-\sin x(\cos x-\sin x)}{(\sin x+\cos x)^2}
=\frac{1}{(\sin x+\cos x)^2}
\]
因此在 \(x=\frac\pi4\) 处的切线斜率为
\(\frac{1}{(\frac{\sqrt2}{2}+\frac{\sqrt2}{2})^2}=\frac12\)，故选 B.
\end{solution}

\begin{problem}
已知函数 \(f(x) = \e^{x} - 1\)，\(g(x) = -x^{2} + 4x - 3\)，若有 \( f(a) = g(b) \)，则 \( b \) 的取值范围为（\quad）

\choices
    {\([2 - \sqrt{2}, 2 + \sqrt{2}]\)}
    {\((2 - \sqrt{2}, 2 + \sqrt{2})\)}
    {\([1,3]\)}
    {\((1,3)\)}
\end{problem}

\begin{answer}
B
\end{answer}

\begin{solution}
因为 \(f(a)=\e^a-1>-1\)，所以 \(g(b)>-1\). 又
\[
g(b)=-b^2+4b-3=1-(b-2)^2
\]
于是 \(1-(b-2)^2>-1\)，即 \((b-2)^2<2\)，故
\(2-\sqrt2<b<2+\sqrt2\). 反之，当
\(2-\sqrt2<b<2+\sqrt2\) 时，\(g(b)>-1\)，取
\(a=\ln\bigl(g(b)+1\bigr)\in\R\)，便有
\(f(a)=\e^{\ln(g(b)+1)}-1=g(b)\). 因而
\(b\) 的取值范围为 \((2-\sqrt2,2+\sqrt2)\)，故选 B.
\end{solution}

\section{选做题}

\begin{problem}
在直角坐标系 \(xOy\) 中，曲线 \(C_1\) 的参数方程为 \(\begin{cases}
x = 2\cos \alpha,\\
y = \sqrt{3}\sin \alpha,
\end{cases}\)（\(\alpha\) 为参数）在极坐标系（与直角坐标系 \(xOy\) 取相同的长度单位，且以原点 \(O\) 为极点，以 \(x\) 轴正半轴为极轴）中，曲线 \(C_2\) 的方程为 \(\rho(\cos \theta - \sin \theta) + 1 = 0\)，则 \(C_1\) 与 \(C_2\) 的交点个数为\fillinblank{}.
\end{problem}

\begin{answer}
2
\end{answer}

\begin{solution}
由极坐标与直角坐标的关系 \(x=\rho\cos\theta\)，\(y=\rho\sin\theta\)，曲线 \(C_2\) 的方程化为
\(x-y+1=0\)，即 \(y=x+1\). 曲线 \(C_1\) 满足
\[
\frac{x^2}{4}+\frac{y^2}{3}=1
\]
代入 \(y=x+1\)，得
\(\frac{x^2}{4}+\frac{(x+1)^2}{3}=1\)，即
\(7x^2+8x-8=0\). 其判别式为
\(8^2-4\times7\times(-8)=288>0\)，所以有两个不同实根，交点个数为 \(2\).
\end{solution}

\begin{problem}
已知某试验范围为 \([10,90]\)，若用分数法进行 \(4\) 次优选试验，则第二次试点可以是\fillinblank{}.
\end{problem}

\begin{answer}
\(40\) 或 \(60\)
\end{answer}

\begin{solution}
试验区间长度为 \(90-10=80\). 设斐波那契数列满足
\(F_0=F_1=1\)，\(F_{n+1}=F_n+F_{n-1}\)，则
\(F_3=3\)，\(F_4=5\)，\(F_5=8\). 分数法进行 \(4\) 次试验时，首轮两个对称试点分别取区间的
\(\frac{F_3}{F_5}=\frac38\) 和 \(\frac{F_4}{F_5}=\frac58\) 处分点；根据试验结果，第二次试点可能是其中任一个. 因此试点可以为
\(10+\frac38\times80=40\)，\(10+\frac58\times80=60\).
所以第二次试点可以是 \(40\) 或 \(60\).
\end{solution}

\section{填空题}

\begin{problem}
若执行如图所示的框图，输入 \(x_{1} = 1\)，\(x_{2} = 2\)，\(x_{3} = 4\)，\(x_{4} = 8\)，则输出的数等于\fillinblank{}.

\bitmapfigure[width=0.30\linewidth]{img/2011/hunan_liberal/q11_fig1.png}
\end{problem}

\begin{answer}
\(\frac{15}{4}\)
\end{answer}

\begin{solution}
框图开始时 \(i=1\)，\(x=0\). 依次执行加法和判断：当 \(i=1\)，\(2\)，\(3\) 时，\(i<4\) 成立，分别把 \(x_1\)，\(x_2\)，\(x_3\) 加入 \(x\) 并令 \(i\) 加 \(1\)；此时 \(i=4\)，再把 \(x_4\) 加入后，条件 \(i<4\) 不成立. 因而
\(x=1+2+4+8=15\)，\(x=\frac{x}{4}=\frac{15}{4}\).
所以输出的数为 \(\frac{15}{4}\).
\end{solution}

\begin{problem}
已知 \( f(x) \) 为奇函数，\( g(x) = f(x) + 9 \)，\( g(-2) = 3 \)，则 \( f(2) = \)\fillinblank{}.
\end{problem}

\begin{answer}
6
\end{answer}

\begin{solution}
由 \(g(-2)=3\) 和 \(g(x)=f(x)+9\)，得
\(f(-2)+9=3\)，所以 \(f(-2)=-6\). 又因为 \(f\) 是奇函数，
\(f(2)=-f(-2)=6\).
\end{solution}

\begin{problem}
设向量 \(\bs{a}\)，\(\bs{b}\) 满足 \(|\bs{a}| = 2\sqrt{5}\)，\(\bs{b} = (2, 1)\)，且 \(\bs{a}\) 与 \(\bs{b}\) 的方向相反，则 \(\bs{a}\) 的坐标为\fillinblank{}.
\end{problem}

\begin{answer}
\((-4,-2)\)
\end{answer}

\begin{solution}
因为 \(\bs{a}\) 与 \(\bs{b}\) 方向相反，存在 \(\lambda>0\)，使
\(\bs{a}=-\lambda\bs{b}\). 由 \(|\bs{b}|=\sqrt{2^2+1^2}=\sqrt5\) 及
\(|\bs{a}|=2\sqrt5\)，得 \(\lambda=2\). 因此
\(\bs{a}=-2(2,1)=(-4,-2)\).
\end{solution}

\begin{problem}
设 \(m > 1\)，在约束条件 \(\begin{cases}
y \geqslant x,\\
y \leqslant mx,\\
x + y \leqslant 1
\end{cases}\) 下目标函数 \(z = x + 5y\) 的最大值为 \(4\)，则 \(m\) 的值为\fillinblank{}.
\end{problem}

\begin{answer}
3
\end{answer}

\begin{solution}
由 \(y\geqslant x\) 和 \(y\leqslant mx\)，且 \(m>1\)，可知可行域中的 \(x\geqslant0\). 其顶点为
\(O(0,0)\)，\(A\left(\frac12,\frac12\right)\)，\(B\left(\frac{1}{m+1},\frac{m}{m+1}\right)\).
在这三个顶点处，目标函数 \(z=x+5y\) 的值分别为
\(0\)，\(3\)，\(\frac{1+5m}{m+1}\).
已知最大值为 \(4\)，而 \(3<4\)，所以
\(\frac{1+5m}{m+1}=4\). 解得 \(1+5m=4m+4\)，即 \(m=3\)，满足 \(m>1\).
\end{solution}

\begin{problem}
已知圆 \(C: x^{2} + y^{2} = 12\)，直线 \(l: 4x + 3y = 25\).
\begin{enumerate}
\item 圆 \(C\) 的圆心到直线 \(l\) 的距离为\fillinblank{}；
\item 圆 \(C\) 上任意一点 \(A\) 到直线 \(l\) 的距离小于 2 的概率为\fillinblank{}.
\end{enumerate}
\end{problem}

\begin{answer}
（1）\(5\)；（2）\(\frac{1}{6}\)
\end{answer}

\begin{solution}
\begin{enumerate}
\item 圆心为 \(O(0,0)\)，所以圆心到直线 \(l\) 的距离为
\[
\frac{|4\times0+3\times0-25|}{\sqrt{4^2+3^2}}=5
\]
\item 圆 \(C\) 的半径为 \(2\sqrt3\). 对圆上点 \(A(x,y)\)，因为
\(4x+3y\leqslant5\cdot2\sqrt3<25\)，故
\[
d(A,l)=\frac{25-4x-3y}{5}
\]
条件 \(d(A,l)<2\) 等价于 \(4x+3y>15\). 令
\(u=(4x+3y)/5\)，则 \(u\) 是半径方向上的坐标，且在圆周上可写为
\(u=2\sqrt3\cos\varphi\). 因而
\[
2\sqrt3\cos\varphi>3
\quad\Longleftrightarrow\quad
\cos\varphi>\frac{\sqrt3}{2}
\quad\Longleftrightarrow\quad
|\varphi|<\frac{\pi}{6}
\]
满足条件的弧长对应圆心角为 \(\pi/3\)，所以所求概率为
\(\frac{\pi/3}{2\pi}=\frac16\).
\end{enumerate}
\end{solution}

\begin{problem}
给定 \(k \in \N^*\)，设函数 \(f: \N^* \to \N^*\) 满足：对于任意大于 \(k\) 的正整数 \(n\)，\(f(n) = n - k\).
\begin{enumerate}
\item 设 \(k = 1\)，则其中一个函数 \(f\) 在 \(n = 1\) 处的函数值为\fillinblank{}；
\item 设 \(k = 4\)，且当 \(n \leqslant 4\) 时，\(2 \leqslant f(n) \leqslant 3\)，则不同的函数 \(f\) 的个数为\fillinblank{}.
\end{enumerate}
\end{problem}

\begin{answer}
（1）\(a\in\N^*\)；（2）\(16\)
\end{answer}

\begin{solution}
\begin{enumerate}
\item 当 \(k=1\) 时，对于所有 \(n>1\)，函数值被确定为 \(f(n)=n-1\). 对 \(n=1\)，题设没有给出额外限制，只要求 \(f(1)\in\N^*\)，所以其中一个函数在 \(n=1\) 处的函数值可以是任意 \(a\in\N^*\).
\item 当 \(k=4\) 时，\(n>4\) 的函数值也都被唯一确定为 \(f(n)=n-4\). 对 \(n=1\)，\(2\)，\(3\)，\(4\)，每个函数值都可独立取 \(2\) 或 \(3\)，共有
\(2^4=16\) 个不同函数.
\end{enumerate}
\end{solution}

\section{解答题}

\begin{problem}
在 \(\triangle ABC\) 中，角 \(A\)，\(B\)，\(C\) 所对的边分别为 \(a\)，\(b\)，\(c\)，且满足 \(c \sin A = a \cos C\).
\begin{enumerate}
\item 求角 \(C\) 的大小；
\item 求 \( \sqrt{3}\sin A - \cos\left(B + \frac{\pi}{4}\right) \) 的最大值，并求取得最大值时角 \(A\)，\(B\) 的大小.
\end{enumerate}
\end{problem}

\begin{solution}
\begin{enumerate}
\item
由正弦定理，存在外接圆半径 \(R\)，使
\(a=2R\sin A\)，\(c=2R\sin C\).
代入条件 \(c\sin A=a\cos C\)，得
\[
2R\sin C\sin A=2R\sin A\cos C
\]
三角形内 \(0<A<\pi\)，故 \(\sin A>0\)，从而
\(\sin C=\cos C\). 又 \(0<C<\pi\)，所以
\[
C=\frac{\pi}{4}
\]

\item
于是 \(A+B=\pi-C=\frac{3\pi}{4}\)，即
\(B=\frac{3\pi}{4}-A\)，且 \(0<A<\frac{3\pi}{4}\). 目标式化为
\[
\sqrt3\sin A-\cos\left(B+\frac{\pi}{4}\right)
=\sqrt3\sin A-\cos(\pi-A)
=\sqrt3\sin A+\cos A
=2\sin\left(A+\frac{\pi}{6}\right)
\]
由于 \(0<A+\frac{\pi}{6}<\frac{11\pi}{12}\)，且
\(A+\frac{\pi}{6}\) 可以取 \(\frac{\pi}{2}\)，目标式最大值为 \(2\). 等号成立时
\(A+\frac{\pi}{6}=\frac{\pi}{2}\)，\(A=\frac{\pi}{3}\)，\(B=\frac{3\pi}{4}-\frac{\pi}{3}=\frac{5\pi}{12}\).
\end{enumerate}
\end{solution}

\begin{problem}
某河流上的一座水力发电站，每年六月份的发电量 \(Y\)（单位：\(\text{万千瓦时}\)）与该河上游在六月份的降雨量 \(X\)（单位：\(\text{毫米}\)）有关，据统计，当 \(X = 70\) 时，\(Y = 460\)；\(X\) 每增加 \(10\)，\(Y\) 增加 \(5\). 已知近 \(20\) 年 \(X\) 的值为：\(140\)，\(110\)，\(160\)，\(70\)，\(200\)，\(160\)，\(140\)，\(160\)，\(220\)，\(200\)，\(110\)，\(160\)，\(160\)，\(200\)，\(140\)，\(110\)，\(160\)，\(220\)，\(140\)，\(160\).

\begin{enumerate}
\item 完成如下的频率分布表：

\begin{center}
\begin{tabular}{c|cccccc}
降雨量 & \(70\) & \(110\) & \(140\) & \(160\) & \(200\) & \(220\) \\
\hline
频率 & \(\frac{1}{20}\) &\fillinblank{} & \(\frac{4}{20}\) &\fillinblank{} &\fillinblank{} & \(\frac{2}{20}\)
\end{tabular}
\end{center}

\item 假定今年六月份的降雨量与近 \(20\) 年六月份降雨量的分布规律相同，并将频率视为概率，求今年六月份该水力发电站的发电量低于 \(490\text{ 万千瓦时}\) 或超过 \(530\text{ 万千瓦时}\) 的概率.
\end{enumerate}
\end{problem}

\begin{solution}
\begin{enumerate}
\item
统计给出的 \(20\) 个数据中，\(110\) 出现 \(3\) 次，\(160\) 出现 \(7\) 次，\(200\) 出现 \(3\) 次，因此频率分布表中缺少的三项依次为
\(\frac{3}{20}\)，\(\frac{7}{20}\)，\(\frac{3}{20}\).

\item
由题意，当 \(X=70\) 时 \(Y=460\)，且 \(X\) 每增加 \(10\)，\(Y\) 增加 \(5\)，所以
\[
Y=460+\frac{5}{10}(X-70)=425+\frac{X}{2}
\]
各个可能的 \(X\) 值对应
\begin{center}
\begin{tabular}{c|cccccc}
\(X\) & \(70\) & \(110\) & \(140\) & \(160\) & \(200\) & \(220\) \\
\hline
\(Y\) & \(460\) & \(480\) & \(495\) & \(505\) & \(525\) & \(535\)
\end{tabular}
\end{center}
条件 \(Y<490\) 或 \(Y>530\) 对应 \(X=70\)，\(110\)，\(220\). 因而所求概率为
\[
\frac{1+3+2}{20}=\frac{6}{20}=\frac{3}{10}
\]
\end{enumerate}
\end{solution}

\begin{problem}
如图，在圆锥 \(PO\) 中，已知 \(PO = \sqrt{2}\)，\(\odot O\) 的直径 \(AB = 2\)，点 \(C\) 在 \(\widehat{AB}\) 上，且 \(\angle CAB = 30^{\circ}\)，\(D\) 为 \(AC\) 的中点.

\begin{enumerate}
\item 证明： \(AC \perp\) 平面 \(POD\)；
\item 求直线 \(OC\) 和平面 \(PAC\) 所成角的正弦值.
\end{enumerate}

\bitmapfigure[width=0.32\linewidth]{img/2011/hunan_liberal/q19_fig1.png}
\end{problem}

\begin{solution}
\begin{enumerate}
\item
在 \(\triangle ABC\) 中，\(AB\) 是直径，所以
\(\angle ACB=90^\circ\). 又 \(\angle CAB=30^\circ\)，故
\[
AC=AB\cos30^\circ=2\cdot\frac{\sqrt3}{2}=\sqrt3
\]
因为 \(D\) 是弦 \(AC\) 的中点，\(O\) 是圆心，所以
\(OD\perp AC\). 又 \(OA=OC\)，且 \(PO\) 垂直于底面，所以
\[
PA^2=PO^2+OA^2=PO^2+OC^2=PC^2
\]
即 \(PA=PC\). 因而 \(\triangle PAC\) 是等腰三角形，\(D\) 为底边
\(AC\) 的中点，故 \(PD\perp AC\). 直线 \(OD\) 与 \(PD\) 相交且都在平面 \(POD\) 内，因此
\[
AC\perp\text{平面 }POD
\]

\item
在直角三角形 \(AOC\) 中，\(OA=OC=1\)，且 \(D\) 是斜边 \(AC\) 的中点，所以
\[
OD=\sqrt{OA^2-AD^2}
=\sqrt{1-\left(\frac{\sqrt3}{2}\right)^2}
=\frac12
\]
又 \(PO\perp OD\)，从而
\[
PD=\sqrt{PO^2+OD^2}
=\sqrt{2+\frac14}
=\frac32
\]
在平面 \(POD\) 内作 \(OH\perp PD\). 由于平面 \(PAC\) 含有
\(AC\)，而 \(AC\perp\) 平面 \(POD\)，所以平面 \(PAC\) 与平面 \(POD\) 的交线
\(PD\) 垂直于 \(AC\)，两平面互相垂直. 因而 \(OH\perp PD\) 时
\(OH\perp\text{平面 }PAC\). 由 \(\triangle OPD\) 的面积可得
\(\frac12\,PO\cdot OD=\frac12\,PD\cdot OH\)，\(OH=\frac{PO\cdot OD}{PD} =\frac{\sqrt2\cdot\frac12}{\frac32} =\frac{\sqrt2}{3}\).
由于 \(C\) 在平面 \(PAC\) 内，\(OC=1\)，直线 \(OC\) 与平面 \(PAC\) 所成角的正弦值为
\[
\frac{\operatorname{dist}(O,\text{平面 }PAC)}{OC}
=\frac{OH}{OC}
=\frac{\sqrt2}{3}
\]
\end{enumerate}
\end{solution}

\begin{problem}
某企业在第\(1\)年初购买一台价值为 \(120\text{ 万元}\) 的设备 \(M\)，\(M\) 的价值在使用过程中逐年减少. 从第\(2\)年到第\(6\)年，每年初 \(M\) 的价值比上年初减少 \(10\text{ 万元}\)；从第\(7\)年开始，每年初 \(M\) 的价值为上年初的 \(75\%\).
\begin{enumerate}
\item 求第 \(n\) 年初 \(M\) 的价值 \(a_{n}\) 的表达式；
\item 设 \(A_{n} = \frac{a_{1} + a_{2} + \cdots + a_{n}}{n}\)，若 \(A_{n}\) 大于 \(80\text{ 万元}\)，则 \(M\) 继续使用，否则须在第\(n\)年初对 \(M\) 更新，证明：须在第\(9\)年初对 \(M\) 更新.
\end{enumerate}
\end{problem}

\begin{solution}
\begin{enumerate}
\item
第 \(1\) 年初的价值为 \(a_1=120\). 从第 \(2\) 年到第 \(6\) 年，每年比上年初少 \(10\) 万元，所以
\(a_n=120-10(n-1)=130-10n\)，\(1\leqslant n\leqslant6\).
第 \(6\) 年初的价值为 \(a_6=70\)，从第 \(7\) 年开始按 \(75\%\) 递减，所以
\(a_n=70\left(\frac34\right)^{n-6}\)，\(n\geqslant7\).
故
\[
a_n=
\begin{cases}
130-10n,&1\leqslant n\leqslant6,\\
70\left(\frac34\right)^{n-6},&n\geqslant7.
\end{cases}
\]

\item
对于 \(1\leqslant n\leqslant6\)，由等差数列求和公式，
\(A_n=\frac{1}{n}\sum_{i=1}^{n}(130-10i) =\frac{130n-5n(n+1)}{n} =125-5n\geqslant95>80\)，\(1\leqslant n\leqslant6\).
随后
\(a_7=\frac{105}{2}\)，\(a_8=\frac{315}{8}\)，\(a_9=\frac{945}{32}\).
因此
\(A_7=\frac{\frac{1245}{2}}7=\frac{1245}{14}>80\)，\(A_8=\frac{\frac{5295}{8}}8=\frac{5295}{64}>80\)
而
\[
A_9=\frac{\frac{22125}{32}}9=\frac{7375}{96}<80
\]
所以第 \(1\) 至第 \(8\) 年初均应继续使用，至第 \(9\) 年初时平均价值不超过 \(80\) 万元，须在第 \(9\) 年初更新设备.
\end{enumerate}
\end{solution}

\begin{problem}
已知平面内一动点 \(P\) 到点 \(F(1,0)\) 的距离与点 \(P\) 到 \(y\) 轴的距离的差等于 \(1\).
\begin{enumerate}
\item 求动点 \(P\) 的轨迹 \(C\) 的方程；
\item 过点 \(F\) 作两条斜率存在且互相垂直的直线 \(l_{1}\)，\(l_{2}\)，设 \(l_{1}\) 与轨迹 \(C\) 相交于点 \(A\)，\(B\)，\(l_{2}\) 与轨迹 \(C\) 相交于点 \(D\)，\(E\)，求 \(\overrightarrow{AD} \cdot \overrightarrow{EB}\) 的最小值.
\end{enumerate}
\end{problem}

\begin{solution}
\begin{enumerate}
\item
设动点 \(P(x,y)\). 由题意
\[
\sqrt{(x-1)^2+y^2}-|x|=1
\]
即
\[
\sqrt{(x-1)^2+y^2}=|x|+1
\]
两边都非负，平方得
\((x-1)^2+y^2=(|x|+1)^2\)，\(y^2=2x+2|x|\).
当 \(x\geqslant0\) 时，\(y^2=4x\)；当 \(x<0\) 时，\(y^2=0\)，即 \(y=0\). 反代可知这两部分都满足原方程，因此轨迹为
\[
C:\quad
\begin{cases}
y^2=4x,&x\geqslant0,\\
y=0,&x<0.
\end{cases}
\]

\item
两条直线均过 \(F(1,0)\)，且斜率存在并互相垂直. 设 \(l_1\) 的斜率为 \(k\)，则 \(k\ne0\)，可写为
\(l_1:y=k(x-1)\)，\(l_2:y=-\frac1k(x-1)\).
负半轴部分 \(y=0\)，\(x<0\) 与任一条过 \(F\) 的直线没有交点；点 \(F\) 本身不在该部分上，所以四个交点均在抛物线上. 令
\[
Q(t)=(t^2,2t)
\]
表示抛物线 \(y^2=4x\) 上的点. 设
\(A=Q(u)\)，\(B=Q(v)\)，\(D=Q(r)\)，\(E=Q(s)\).
由两条直线的方程，\(u\)，\(v\) 是
\[
kt^2-2t-k=0
\]
的两根，\(r\)，\(s\) 是
\[
t^2+2kt-1=0
\]
的两根，因而
\(u+v=\frac2k\)，\(uv=-1\)，\(r+s=-2k\)，\(rs=-1\).
又
\(\overrightarrow{AD}=(r-u)(r+u,2)\)，\(\overrightarrow{EB}=(v-s)(v+s,2)\).
因此
\[
\begin{aligned}
\overrightarrow{AD}\cdot\overrightarrow{EB}
&=(r-u)(v-s)\bigl((r+u)(v+s)+4\bigr)\\
&=(rv+us+2)^2,
\end{aligned}
\]
其中使用了 \(uv=rs=-1\). 记
\(X=rv+us\)，\(Y=ru+vs\).
则
\(X+Y=(r+s)(u+v)=-4\)，\(X+2=\frac{X-Y}{2}\).
此外
\[
\begin{aligned}
(X-Y)^2
&=(r-s)^2(v-u)^2\\
&=\bigl((r+s)^2-4rs\bigr)\bigl((u+v)^2-4uv\bigr)\\
&=16\frac{(k^2+1)^2}{k^2}.
\end{aligned}
\]
所以
\[
\overrightarrow{AD}\cdot\overrightarrow{EB}
=(X+2)^2
=4\frac{(k^2+1)^2}{k^2}
=4\left(k+\frac1k\right)^2
\]
由于 \(k\ne0\)，
\[
\left(k+\frac1k\right)^2-4
=\frac{(k^2-1)^2}{k^2}\geqslant0
\]
故内积的最小值为 \(16\)，当且仅当 \(k=\pm1\) 时取得.
\end{enumerate}
\end{solution}

\begin{problem}
设函数 \( f(x) = x - \frac{1}{x} - a \ln x  \)（\(a \in \R\)）.
\begin{enumerate}
\item 讨论函数 \( f(x) \) 的单调性.
\item 若 \( f(x) \) 有两个极值点 \( x_{1} \) 和 \( x_{2} \)，记过点 \( A(x_{1}, f(x_{1})) \)，\( B(x_{2}, f(x_{2})) \) 的直线斜率为 \( k \). 问：是否存在 \( a \)，使得 \( k = 2 - a \)？若存在，求出 \( a \) 的值；若不存在，请说明理由.
\end{enumerate}
\end{problem}

\begin{solution}
\begin{enumerate}
\item
函数定义域为 \(x>0\)，且
\[
f'(x)=1+\frac1{x^2}-\frac ax
=\frac{x^2-ax+1}{x^2}
\]
因为分母 \(x^2>0\)，只需讨论
\(x^2-ax+1\) 的符号.

当 \(a\leqslant2\) 时，由
\(x^2-ax+1=x\left(x+\frac1x-a\right)\)，\(x+\frac1x\geqslant2\)
可知 \(x^2-ax+1\geqslant0\). 当 \(a<2\) 时导数处处为正；当 \(a=2\) 时仅在 \(x=1\) 处导数为 \(0\)，函数在 \((0,+\infty)\) 上仍为增函数.

当 \(a>2\) 时，方程 \(x^2-ax+1=0\) 的两个根为
\(x_1=\frac{a-\sqrt{a^2-4}}2\)，\(x_2=\frac{a+\sqrt{a^2-4}}2\)
满足 \(0<x_1<x_2\)，且 \(x_1x_2=1\). 导数符号依次为正、负、正，所以
\begin{center}
\begin{tabular}{c|ccc}
\(x\) & \((0,x_1)\) & \((x_1,x_2)\) & \((x_2,+\infty)\) \\
\hline
\(f'(x)\) & + & - & +
\end{tabular}
\end{center}
函数在 \((0,x_1)\) 和 \((x_2,+\infty)\) 上递增，在 \((x_1,x_2)\) 上递减.

\item
若 \(f\) 有两个极值点，则必有 \(a>2\)，并令
\(x_1=\e^{-t}\)，\(x_2=\e^t\)，\(t>0\).
于是 \(a=x_1+x_2=2\cosh t\). 由 \(x_1x_2=1\)，
\[
\begin{aligned}
k
&=\frac{f(x_2)-f(x_1)}{x_2-x_1}\\
&=2-a\frac{\ln(x_2/x_1)}{x_2-x_1}\\
&=2-a\frac{t}{\sinh t}.
\end{aligned}
\]
因为 \(t>0\) 时 \(\sinh t>t\)，所以
\[
k>2-a
\]
因此不可能有 \(k=2-a\)，不存在满足条件的 \(a\).
\end{enumerate}
\end{solution}
